\documentclass[twocolumn,showpacs,amsmath,amssymb]{revtex4}
%\documentclass[preprint,showpacs,preprintnumbers,amsmath,amssymb]{revtex4}
\usepackage{graphicx}
%\usepackage{rotating}
\begin{document}
\title {Analytic scattering model with many closed channels}
\author{Christopher Ticknor}
\affiliation{Theoretical Division, Los Alamos National Laboratory, Los Alamos, New Mexico 87545, USA}
\affiliation{Kavli Institute of Theoretical Physics, University of California,
Santa Barbara, Santa Barbara, California 93106, USA}
\date{\today}
\begin{abstract}
We present results from an N channel finite square well scattering model 
where only one channel is open.
By using spatially independent couplings are short range, 
the scattering is nearly analytic and the density of states can be tuned.
This allows for direct comparison between RMT and scattering results.  
Such models has important implications for molecular scattering. 
\end{abstract}
%\pacs{03.75Hh,67.85-d}
\maketitle
%%%%%%%%%%%%%%%%%%%%%%%%%%%%%%
\section{Introduction}
I have tried to say everything that is needed to setup the scattering problem. 
It is rather dense...

Study simple model and see what happens when one
has many closed channels.
This model can be used gain an understanding of highly resonant scattering.

One more important addition is the two length scale model.
This will strongly mix the bound state's character.

Then compare to RMT theory.   

\section{The model}
This is a simple N square well model with 
$N$ channels only one of which is open.
This is a generalization of the 2 channel square well from
Ref.  [PHYSICAL REVIEW A,  65, 053617].

An interesting feature to this model is that there is an eigenbasis
at short range. What this means is that the DOS/bound states will 
not be strongly coupled. Their "coupling" only occurs at $R$ where 
they are closed. So they spectrum this model generates, might not
have the properties we want. A way to over come this is to make
the model a multi-step square well, or add another length scale to the problem. 
This way there will another step that totally scrambles the states and therefore should make the DOS or bound states spectrum strongly coupled. 

It is composed of kinetic and potential energy terms
\begin{eqnarray}
[T_{ij}+V_{ij}(r)]\psi_j(r)=E \psi_{i}(r)
\end{eqnarray}
Where $T_{ij}=-{\hbar^2 \over2m} {\nabla^2}\delta_{ij}$
is the kinetic energy operator and $i=\{c,o\}$ denotes the set of closed channels and the open channel.
$T$ is diagonal in open/closed channel space, this enforced with Kronecker delta. 
We will pick the units so that $\hbar$ is 1
and $m$ is set to 1 (this can be changed easily). 

The parameters that define the model are
$R$, the width of the potential well;
$C_{ij}$, the coupling between channels;
$D_i$, the depths of the potential wells;
$\Delta_i$, the closed threshold energy; and
$E$, the scattering energy and wavenumber $k_o=\sqrt{2mE}$.
We are able to tune the density of states with good chooses of
$R$, $C_{ij}$ and $D_j$.

The interaction for $r<R$ is
\begin{eqnarray}
V_{ii}(r)=-D_i
\\\nonumber
V_{ij}(r)=C_{ij}.
\end{eqnarray}
For $r>R$, it is
\begin{eqnarray}
V_{ii}(r)=\Delta_i\delta_{ic}
\\\nonumber
V_{ij}=0.
\end{eqnarray}
Here $\delta_{ic}$ is 1 when $i$
is a closed channel and zero otherwise. 
For all cased considered $C_{ij}$ is symmetric and real.

We can pick the orbital angular momentum of this problem and add a centrifugal 
barrier simply by changing the matching functions. This leads to another
parameters: $L$.  The barrier height is $L(L+1)/2mR^2$ and this can be tuned.

The interaction matrix, $V_{ij}$, is an $N\times N$ matrix, which is diagonalized. 
From this we get the eigenvalues: $\epsilon_\alpha$.  Furthermore, 
the eigenvectors can be used to form a frame transformation matrix, $U_{i\alpha}$, that defines how one transforms between the eigenbasis ($\alpha$) to the open/closed basis ($i$).

The wavefunctions for $r<R$ are Riccati-Bessel functions ($f_\alpha$) with
wave number $k_\alpha=\sqrt{2m\epsilon_\alpha}$.
Then the wavefunction  is $\phi_{i\alpha}(r)=U_{i\alpha}f_\alpha(r)$ for the $\alpha$ solution.

First we will setup the general equations before reformatting them to be solved.
To correctly solve the full problem, we must take a linear combination of
the inside wavefunctions so that they smoothly match to both the decaying 
wavefunctions outside of $R$ 
and have the desired scattering boundary conditions for the open channel.
For $L=0$, the equations we need to solve for a closed channel ($c$) are 
$\sum_\alpha b_\alpha \phi_{c\alpha}(R)=a_c e^{-k_c R}$ and
$\sum_\alpha b_\alpha \phi_{c\alpha}^\prime(R)=- a_c k_c e^{-k_c R}$
where $k_c=\sqrt{2m(\Delta_c-E)}$. 
Here $b_\alpha$ are the coefficients denoting the linear combination of
the inside wavefunctions
and $a_c$ is the amplitude of the closed channel.

In matrix form, the equations are:
\begin{eqnarray}
\Phi_{m\alpha} b_\alpha= \Psi_{mj} a_j.
\end{eqnarray}
$\Phi_{m\alpha}$ is $2N\times N$ matrix made up of 
all $N$ channels plus $N$ radial derivates for each channel.
$\Phi_{m\alpha}$ is made up of 
$\phi_{i\alpha}(R)$ and $\phi^\prime_{i\alpha}(R)$.
$m$ is an index the encompasses both $i$
and the radial derivative and therefore index $m$ ranges 1 and $2N$.  
$\alpha$ covers all $N$ eigenstates.
So a column of $\Phi_{m\alpha}$ is
$\{\phi_{1\alpha}(R),\phi^\prime_{1\alpha}(R),
\phi_{2\alpha}(R),\phi^\prime_{2\alpha}(R),...,
\phi_{N\alpha}(R),\phi^\prime_{N\alpha}(R)\}^T$.

$\Psi_{mj}$ is a matrix the contains the correct physics behavior 
outside of $R$. It is an $2N\times (N+1)$ matrix for an $N$ channel problem with one open channel.
There is an extra coefficient for the open channel scattering boundary condition. 
For the closed channels, $\Psi_{mc}=g_c$ 
where $g_c$ is the decaying function for $c^{th}$ closed channel.
For the open channel, we have
$\Psi_{mo}=c f_o- s g_o$ and $\Psi_{mo+1}^\prime =c f_o^\prime - s g_o^\prime$
where $f_o$ ($g_o$) is the regular (irregular) scattering function that are
the Riccati-Bessel functions.
$c$ and $s$ combine to make $s/c=\tan(\delta)$ where $\delta$ is the scattering phase shift. So a column of $\Psi_{m1}$ for the first closed channel is
$(g_{1}(R),g_1^\prime(R),0,..., 0,0)^T$.
For the open channel, we have
 $\Psi_{m,N}$ is
 $\{0,0,...,0,f_{o}(R),f_o^\prime(R)\}^T$
and $\Psi_{m,N+1}$ is
$\{0,0,...,0,g_{o}(R),g_o^\prime(R)\}^T$.
Note, we have picked the open channel to the $N^{th}$ channel.

%For the general problem of more open channels one must.
%We will use the 

The unknowns are $b_\alpha$ and $a_j$.
$b_\alpha$ are the coefficients that tell us how to combine 
the inside wavefucntions to solve the desired boundary conditions
specified by $\Psi_j$.
$a_j$ contains $N-1$ closed channels ($a_c$) amplitudes plus
the scattering  coefficients ($c$ and $s$) for a total of $N+1$ unknowns. 
To solve this problem, we must add one more equation so that
we have $2N+1$ equations for the $2N+1$ unknowns: $b_\alpha$
and $a_j$ (=$\{a_c, s,c\}$).  We typically take $a_1=1$, but 
this choice is arbitrary. Alternately, one could just let $c=1$
and $s$ would become the $\tan(\delta)$.

This can now be set up as standard $Ax=b$ linear solve. 
Picking one coefficient to be one, makes the $b$ vector to have 
one element that is 1 and the rest are zero.
$A$ is $2N+1$ by $2N+1$ matrix.
It is composed of $\Phi_{m\alpha}$ for the upper left block of $A_{}$.
The upper right block is $-\Psi_{mj}$, and this forms all but the last row of the matrix.  
The last row of $A$ is simply all zero except for one element set to one, e.g. picking $a_1$=1 and our indexing scheme, we have $A_{2N+1,N+1}=1$.
So now one solves for $x$, it is $\{b_\alpha,a_c,c,s\}^T$.

After solving the problem, we can renormalize the final wavefunction
so that the scattering amplitude is 1 or energy normalized?
$c^2+s^2=1$.
Energy normalization is done to reflect the density of states of the scattering 
wavefunctions. 

With our notation the full wavefunction can be written as: for $r<R$
\begin{eqnarray}
\psi_i(r)=\sum_\alpha b_\alpha \phi_{i\alpha}(r)
\end{eqnarray}
here $i=\{c,o\}$ and $\phi_{i\alpha}=\sum_\beta U_{i\beta}[\delta_{\beta\alpha}f_\alpha(r)]$.
For $r>R$ the wavefunction  is:
\begin{eqnarray}
&\psi_{c}(r)=a_c g_c(r)
\\\nonumber
&\psi_o(r)=c f_o(r)-s g_o(r)  
\end{eqnarray}
Here $c$ is a index for all closed channels.
 
One important measure is the closed channel amplitude:
\begin{eqnarray}
Z=\sum_c \int_0^\infty  |\psi_c(r)|^2 dr.
\end{eqnarray}
The scattering cross section is 
\begin{eqnarray}
\sigma={4 \pi \over k^2} \sin^2(\delta)
\end{eqnarray}

We now discuss some standard cases we study in this work for the 
short range interaction.

{\it Tri-diagonal case}:
we take all $D_i$ be the same for all channels ($D$).
For the simplest case, we take $C_{ij}$
to be constant ($C$) and non-zero when $j=i\pm1$,
where $j$and $i$ are the channel indices.  
For the long range interaction we need the scattering energy $\epsilon$ and $\Delta_i$ (also taken to be constant for all channels).

{\it GOE case}: $C_{ij}$ is defined by gaussian orthogonal ensemble. 
It is a bunch of random numbers, but real and symmetric.
We usually take all channels to have one depth as well ($D$). 

{\it Perturbatively coupled open/closed case}:
This is where we solve a general closed channel system and 
In this case there is a general $C_{ij}$ matrix, and for the closed-closed sub-block
we can solve the complete bound state problem with $N-1$ closed channels.
We can then perturbatively couple the known bound, closed channel system to 
the open channel by adding on $\lambda C_{oc}$ where $o$ and $c$ denote all
the open and closed channels and $\lambda$ is a parameter that can be tuned 
form 0 to 1. 
This allows us to tune a known density of resonances explicitly. 
It also then allows us to see how coupling to an open channel influences 
multi-channel molecular bound states.

\section{Two length scale}
we can introduce another radius $S$ and it has the same basic
properties of the first well, but they are different. So like different depths
and off-diagonal channel couplings. 
Now the eigenbasis is given by matrix $V$ (instead of $U$).
This scrambles the quantum number numbers and mixes the states
furthers. much like rotational mixing or some other longer range interaction. 

To solve this problem one would simply take the previous solution
at short range $\Phi_\alpha(R)$ and use the new basis transformation
to project this the wavefunction onto the new eigenfunctions at $R$.
This projection then establishes the boundary condition
for each channel (and each $\alpha$) to match to the  
new basis functions. Then we can a propagate the wavefucntion
in each channel to $S$. At this point we recover the exact same problem
that we had before, but with a new $\Phi_{m\alpha}$.


\section{DOS estimates}
 
Once we have the eigenfunctions in the inner region, we are able to find
the bound-states if we assume all channels terminate at $R$ (infinite square well). Then we are able easily estimate the resonances.
We can solve the full problem by correctly searching....
While not correct, modifications are only slight.
it is reasonable for large $\Delta$.
It offers a very good estimate of density of states.


\section{Energy dependence}

\begin{figure}%{r}{0.45\textwidth}
%\includegraphics[width=0.95\columnwidth]{sr0_field.png}
\caption{The cross section}
\label{energy1}
\end{figure}

\bibliography{/Users/cticknor/BIB/coldBIB}{}
\bibliographystyle{apsrev.bst}
\end{document}
